\documentclass[10pt,a4paper]{amsart}
\usepackage[margin=19mm]{geometry}
\usepackage{amsmath,amssymb,mathtools}
\usepackage[scr=rsfs,cal=euler]{mathalfa}
\usepackage{xfrac}
\usepackage[unicode,hidelinks,bookmarks=false]{hyperref}
\newcommand{\ie}{i.e.\ }
\newcommand{\cf}{cf.\ }
\newcommand{\resp}{respectively\ }
\newcommand{\Subsection}{\subsection}
\providecommand{\nobreakdash}{\nobreak\hbox{-}}
\setlength{\emergencystretch}{2em}
\multlinegap0pt
\usepackage[shortlabels]{enumitem}
\usepackage{tikz}
\usepackage{amssymb}
\usepackage{xfrac}
\usepackage{dsfont}
\usepackage{mathtools}
\newtheorem{lemma}{Lemma}
\newtheorem{proposition}{Proposition}
\DeclareMathOperator{\length}{length}
\DeclareMathOperator{\mult}{mult}
\title{The Milnor Number of One Dimensional Local Rings}
\author{Yotam Svoray}
\date{Source: arXiv:2602.11519v1 (CC BY 4.0)}
\begin{document}
\maketitle

\noindent\emph{Source and licence.} The Milnor Number of One Dimensional Local Rings. Version arXiv:2602.11519v1 (CC BY 4.0), distributed by arXiv under the Creative Commons Attribution 4.0 International licence, http://creativecommons.org/licenses/by/4.0/, as recorded in the arXiv licence metadata for this exact version and shown on the abstract page https://arxiv.org/abs/2602.11519v1. Full licence notice: \texttt{licenses/arxiv-2602.11519-CC-BY-4.0.txt}.\par\smallskip

\noindent\textit{Editorial scope.} Counted unit: the article's proposition thm:branch (source0.tex lines 281--287). The article's complete original proof of it is delivered with it (source0.tex lines 290--312, one \texttt{proof} environment of 23 lines). No mathematics is written, repaired or re-derived: the statement and the proof are the article's own lines, in the article's own notation, and no retained line is substituted or re-flowed.  Retained uncounted as the source-local context the proof needs: lines 179--181.  Nothing was omitted from any retained span, so the package carries no omission record.  No retained line contains a citation, so the wrapper carries no bibliography and no bibliography entry.  No retained line contains a figure environment, an image-inclusion command or drawing code, so no image is delivered and the package declares no asset dependency.  Editorial formatting only, in addition: an AMS \texttt{amsart} preamble that restates the article's own declarations and macros used by the retained lines, namely the environments \texttt{lemma}, \texttt{proposition} on separate counters, together with the standard packages those lines need.  The retained environments print in this document as Lemma 1, Proposition 1 in order of appearance, whereas the article prints the same items with the numbering of its own full document.  The complete native source of the article as distributed in the arXiv e-print for this exact version is bundled unchanged as \texttt{sources/194478/source0.tex}.
\begin{lemma}\label{lem:mult_sum_prim}
    If $\mathfrak{p}_1, \dots, \mathfrak{p}_{r(R)}$ are the minimal primes of $R$, then we have that $\mult(R)=\sum_{i=1}^r \mult\left(\frac{R}{\mathfrak{p}_i}\right)$. 
\end{lemma}

\begin{proposition}[Hironaka's Lemma] \label{thm:branch}
    Let $\mathfrak{p}_1, \dots, \mathfrak{p}_{r(R)}$ be the minimal primes of $R$. Then we have that 
    \begin{equation*}
        \delta(R) = \sum_{i=1}^{r(R)}  \delta\left(\frac{R}{\mathfrak{p}_i}\right) + \sum_{i <  j} i(\mathfrak{p}_i, \mathfrak{p}_j).
    \end{equation*}
    In addition, $\delta(R) \geq \frac{r(R)(r(R)-1)}{2}$. 
\end{proposition}

\begin{proof}
As in Lemma~\ref{lem:mult_sum_prim}, denote $\mathfrak{q}=\bigcap_{i=1}^{r(R)-1} \mathfrak{p}_i$. Then we that $\mathfrak{q}\cap \mathfrak{p}_{r(R)} = \{0\}$, and so we get a short exact sequence:
\begin{equation*}
    0 \to R \to \frac{R}{\mathfrak{q}} \oplus \frac{R}{\mathfrak{p}_{r(R)}} \to \frac{R}{\mathfrak{q}+\mathfrak{p}_{r(R)}} \to 0. 
\end{equation*}
\noindent Therefore, we can conclude that 
\begin{equation*}
    \delta(R)=\length_R\left(\frac{\overline{R}}{R}\right) = \length_R\left(\frac{\overline{R}}{\frac{R}{\mathfrak{q}} \oplus \frac{R}{\mathfrak{p}_{r(R)}} }\right) + \length_R \left(\frac{\frac{R}{\mathfrak{q}} \oplus \frac{R}{\mathfrak{p}_{r(R)}} }{R} \right).
\end{equation*}

\noindent Yet, the normalization of $\frac{R}{\mathfrak{q}} \oplus \frac{R}{\mathfrak{p}_{r(R)}}$ is exactly $\overline{R}$, that is, $\overline{R} = \overline{\left(\frac{R}{\mathfrak{q}}\right)} \times \overline{\left(\frac{R}{\mathfrak{p}_{r(R)}}\right)}$. In addition, we have that


\begin{equation*}
    \frac{\frac{R}{\mathfrak{q}} \oplus \frac{R}{\mathfrak{p}_{r(R)}} }{R}  = \frac{R}{\mathfrak{q} + \mathfrak{q}_{r(R)}}. 
\end{equation*}
\noindent Therefore, we can conclude that 

\begin{equation*}
    \delta(R) = \delta\left(\frac{R}{\mathfrak{q}}\right) + \delta\left(\frac{R}{\mathfrak{p}_{r(R)}} \right) + i_R(\mathfrak{p}_{r(R)}, \mathfrak{q}).
\end{equation*}
\noindent Since the minimal primes of $\frac{R}{\mathfrak{p}_{r(R)}}$ correspond to $\mathfrak{p}_1, \dots, \mathfrak{p}_{r(R)-1}$, we repeat this process for $\frac{R}{\mathfrak{p}_{r(R)}}$, and the result would follow. For the second part, observe that $i(\mathfrak{p}_i, \mathfrak{p}_j)>0$ for every $i  \neq j$, and so $\sum_{i <  j} i(\mathfrak{p}_i, \mathfrak{p}_j) \geq \frac{r(R)(r(R)-1)}{2}$. 
\end{proof}

\end{document}
